\chapter{Spinal Decompression: Laminectomy and Discectomy}
\section*{Overview}
Spinal decompression removes bone, ligament, or disc material that compresses the spinal cord or nerve roots.
\section*{Common Instruments and Equipment}
The exact setup varies by technique and setting. Common items may include:
\begin{itemize}[nosep,leftmargin=*]
  \item spinal retractor and bone-cutting instruments
  \item microscope or loupe magnification when used
  \item fluoroscopy or navigation equipment as indicated
\end{itemize}

\section*{Why It Is Done}
It may be considered for spinal stenosis, herniated disc with persistent neurologic symptoms, cauda equina syndrome, or progressive weakness after appropriate evaluation.
\section*{How It Is Performed}
Through an open or minimally invasive approach, the surgeon removes part of the lamina and/or herniated disc material. Fusion may be added when instability is present.
\section*{Contraindications}
No single contraindication applies to every patient; the treating team assesses whether the procedure is suitable for the individual and the reason for treatment.

\section*{Precautions}
The care team reviews relevant health conditions, medicines, allergies, and procedure-specific risks, and explains any changes needed for the individual.

\section*{Risks and Recovery}
Dural tear, nerve injury, infection, bleeding, recurrent disc herniation, spinal instability, and persistent symptoms are possible. Mobilization and rehabilitation are tailored to the operation.
\section*{References}
\begin{enumerate}\item MedlinePlus. Laminectomy. U.S. National Library of Medicine; 2025.\item Current Medical Diagnosis and Treatment. McGraw-Hill; 2025.\end{enumerate}
\clearpage
